% Options for packages loaded elsewhere
\PassOptionsToPackage{unicode}{hyperref}
\PassOptionsToPackage{hyphens}{url}
\documentclass[
  ignorenonframetext,
]{beamer}
\newif\ifbibliography
\usepackage{pgfpages}
\setbeamertemplate{caption}[numbered]
\setbeamertemplate{caption label separator}{: }
\setbeamercolor{caption name}{fg=normal text.fg}
\beamertemplatenavigationsymbolsempty
% remove section numbering
\setbeamertemplate{part page}{
  \centering
  \begin{beamercolorbox}[sep=16pt,center]{part title}
    \usebeamerfont{part title}\insertpart\par
  \end{beamercolorbox}
}
\setbeamertemplate{section page}{
  \centering
  \begin{beamercolorbox}[sep=12pt,center]{section title}
    \usebeamerfont{section title}\insertsection\par
  \end{beamercolorbox}
}
\setbeamertemplate{subsection page}{
  \centering
  \begin{beamercolorbox}[sep=8pt,center]{subsection title}
    \usebeamerfont{subsection title}\insertsubsection\par
  \end{beamercolorbox}
}
% Prevent slide breaks in the middle of a paragraph
\widowpenalties 1 10000
\raggedbottom
\AtBeginPart{
  \frame{\partpage}
}
\AtBeginSection{
  \ifbibliography
  \else
    \frame{\sectionpage}
  \fi
}
\AtBeginSubsection{
  \frame{\subsectionpage}
}
\usepackage{iftex}
\ifPDFTeX
  \usepackage[T1]{fontenc}
  \usepackage[utf8]{inputenc}
  \usepackage{textcomp} % provide euro and other symbols
\else % if luatex or xetex
  \usepackage{unicode-math} % this also loads fontspec
  \defaultfontfeatures{Scale=MatchLowercase}
  \defaultfontfeatures[\rmfamily]{Ligatures=TeX,Scale=1}
\fi
\usepackage{lmodern}
\ifPDFTeX\else
  % xetex/luatex font selection
\fi
% Use upquote if available, for straight quotes in verbatim environments
\IfFileExists{upquote.sty}{\usepackage{upquote}}{}
\IfFileExists{microtype.sty}{% use microtype if available
  \usepackage[]{microtype}
  \UseMicrotypeSet[protrusion]{basicmath} % disable protrusion for tt fonts
}{}
\makeatletter
\@ifundefined{KOMAClassName}{% if non-KOMA class
  \IfFileExists{parskip.sty}{%
    \usepackage{parskip}
  }{% else
    \setlength{\parindent}{0pt}
    \setlength{\parskip}{6pt plus 2pt minus 1pt}}
}{% if KOMA class
  \KOMAoptions{parskip=half}}
\makeatother
\setlength{\emergencystretch}{3em} % prevent overfull lines
\providecommand{\tightlist}{%
  \setlength{\itemsep}{0pt}\setlength{\parskip}{0pt}}
\usepackage{bookmark}
\IfFileExists{xurl.sty}{\usepackage{xurl}}{} % add URL line breaks if available
\urlstyle{same}
\hypersetup{
  hidelinks,
  pdfcreator={LaTeX via pandoc}}

\author{\texorpdfstring{}{}}
\date{}

\begin{document}

\section{Conclusions}\label{sec:conclusions}

\begin{frame}{Summary of Contributions}
\protect\phantomsection\label{summary-of-contributions}
This paper introduced a systematic \(2 \times 2\) matrix framework for
analyzing Active Inference's meta-level operation, organized along
Data/Meta-Data and Cognitive/Meta-Cognitive dimensions. The four
resulting quadrants provide a comprehensive taxonomy:

\begin{enumerate}
\item
  \textbf{Quadrant 1 (Data-Cognitive):} Baseline EFE computation with
  direct sensory processing, establishing the foundation for
  pragmatic-epistemic balance through standard policy selection.
\item
  \textbf{Quadrant 2 (Meta-Data-Cognitive):} Extended EFE with meta-data
  weighting and quality integration, demonstrating that incorporating
  confidence scores and reliability metrics improves decision accuracy
  from 85\% to 94\% under uncertain conditions.
\item
  \textbf{Quadrant 3 (Data-Meta-Cognitive):} Hierarchical EFE with
  self-assessment and adaptive control, enabling agents to monitor
  inference quality and switch strategies when confidence degrades.
\item
  \textbf{Quadrant 4 (Meta-Data-Meta-Cognitive):} Framework-level
  optimization enabling cognitive architecture evolution, yielding a
  23\% overall performance improvement through recursive self-analysis
  of meta-cognitive parameters.
\end{enumerate}

Beyond the quadrant taxonomy, this work established three principal
results. First, that Active Inference is meta-pragmatic---the
specification of matrix \(C\) defines complex value landscapes that
transcend scalar reward functions, making what matters a design variable
rather than an external constraint. Second, that Active Inference is
meta-epistemic---the specification of matrices \(A\), \(B\), and \(D\)
determines epistemological boundaries, making what can be known an
explicit modeling choice. Third, that these meta-level properties have
direct implications for cognitive security and AI safety, with each
quadrant presenting distinct attack surfaces and defense requirements.
\end{frame}

\begin{frame}{Computational Validation}
\protect\phantomsection\label{computational-validation}
The framework was validated through implementations yielding
statistically significant results across all hypotheses: meta-data
integration improves performance (Cohen's \(d = 1.05\)), meta-cognitive
control enhances robustness (\(F(3,96) = 12.45\), \(p < 0.001\)), and
framework optimization provides adaptive advantage (Cohen's
\(d = 0.85\)). The performance regression model achieves \(R^2 = 0.87\),
confirming that meta-level processing accounts for substantial variance
in cognitive performance.
\end{frame}

\begin{frame}{Unified Framework}
\protect\phantomsection\label{unified-framework}
Active Inference, through its meta-level operation, provides a unified
framework for understanding:

\begin{itemize}
\tightlist
\item
  \textbf{Perception as Inference:} Bayesian hypothesis testing under
  generative model constraints
\item
  \textbf{Action as Free Energy Minimization:} Goal-directed behavior
  shaped by preference priors
\item
  \textbf{Learning as Model Refinement:} Generative model adaptation
  through experience
\item
  \textbf{Meta-Cognition as Self-Modeling:} Recursive cognitive
  awareness enabling adaptive control
\end{itemize}
\end{frame}

\begin{frame}{Broader Vision}
\protect\phantomsection\label{broader-vision}
The capacity to specify epistemic frameworks (what can be known) and
pragmatic landscapes (what matters) elevates Active Inference from a
theory of cognition to a \emph{meta-theory}---a methodology for
understanding how cognitive theories themselves are constructed,
evaluated, and revised. This meta-theoretical status carries practical
weight: for AI safety, it grounds value alignment in transparent,
inspectable framework parameters rather than opaque reward signals; for
cognitive security, it reveals attack surfaces and defense strategies
that operate at the level of cognitive architecture rather than
individual beliefs.

Ultimately, the quadrant framework developed here suggests that
intelligence is best understood as \emph{framework flexibility}: the
capacity to modify the structures within which cognition operates. As
Active Inference matures as both a theoretical and computational
paradigm, the meta-level perspective offered here provides a foundation
for understanding not just how agents think, but how the conditions of
thought itself are specified, secured, and transformed. Realizing this
promise will require the empirical, computational, and
cross-disciplinary advances outlined in Section
\ref{sec:future_directions}---work that we hope this framework helps to
motivate and organize.
\end{frame}

\end{document}
